On définit une fonction $T$ sur les entiers naturels non nuls par :
\[
T(n) =
\begin{cases}
\dfrac{n}{2}, & \text{si } n \text{ est pair,}\\[4pt]
3n + 1, & \text{si } n \text{ est impair.}
\end{cases}
\]
La \textbf{conjecture de Collatz} affirme que, pour tout entier
$n \geqslant 1$, si l'on applique $T$ de façon répétée à partir de $n$, on
finit toujours par atteindre $1$ ; la suite continue alors dans le cycle
\[
1 \ \to\ 4 \ \to\ 2 \ \to\ 1 \ \to\ \cdots
\]
\begin{enumerate}
  \item Calculer la suite obtenue à partir de $n = 6$, puis de $n = 7$.
  Combien d'applications de $T$ faut-il pour atteindre $1$ ?
  \item Démontrer que, si $n = 2^k$ avec $k \in \mathbb{N}$, on atteint $1$
  après exactement $k$ étapes.
  \item Démontrer que, si $n$ est impair, $T(n)$ est pair, et en déduire
  $T\bigl(T(n)\bigr)$ en fonction de $n$.
  \item \emph{Algorithmique.} Écrire une fonction Python qui renvoie le
  nombre d'étapes nécessaires pour atteindre $1$ à partir de $n$. L'utiliser
  pour $n = 27$ et déterminer la plus grande valeur atteinte.
  \item La conjecture a été vérifiée par ordinateur pour des milliards de
  milliards d'entiers. Est-elle pour autant démontrée ?
\end{enumerate}
